\documentclass[12pt,letterpaper]{article}
\usepackage[utf8]{inputenc}
\usepackage[T1]{fontenc}
\usepackage[spanish]{babel}
\usepackage[margin=3cm]{geometry}
\usepackage{graphicx}
\usepackage{enumitem}
\usepackage{amsmath,amssymb}
\usepackage{array}
\usepackage{tikz}
\usetikzlibrary{arrows.meta,positioning,calc}
\usepackage{listings}
\usepackage[colorlinks=true,linkcolor=black,urlcolor=blue,citecolor=black]{hyperref}
\usepackage{xurl}
\usepackage{titlesec}
\usepackage{parskip}
\usepackage{xcolor}
\titleformat{\section}{\large\bfseries}{\thesection.}{0.5em}{}
\titleformat{\subsection}{\normalsize\bfseries}{\thesubsection.}{0.5em}{}
\emergencystretch=2em
\setlist{itemsep=2pt,topsep=3pt}
\lstset{
  language=C, basicstyle=\ttfamily\footnotesize, keywordstyle=\color{blue!70!black},
  commentstyle=\color{green!40!black}, stringstyle=\color{red!60!black},
  frame=single, showstringspaces=false, breaklines=true, columns=fullflexible, keepspaces=true,
  xleftmargin=0.5em, xrightmargin=0.5em, aboveskip=6pt, belowskip=4pt,
  numbers=left, numberstyle=\tiny\color{gray}, numbersep=6pt,
  literate={á}{{\'a}}1 {é}{{\'e}}1 {í}{{\'i}}1 {ó}{{\'o}}1 {ú}{{\'u}}1
           {Á}{{\'A}}1 {É}{{\'E}}1 {Í}{{\'I}}1 {Ó}{{\'O}}1 {Ú}{{\'U}}1 {ñ}{{\~n}}1 {¿}{{?`}}1,
}
\newcommand{\cod}[1]{\texttt{#1}}

\begin{document}

\begin{titlepage}
\centering
\includegraphics[width=4cm]{escudo.png}\\[0.6cm]
{\Large\bfseries Universidad de Sonora}\\[0.3cm]
{\large Licenciatura en Ciencias de la Computación}\\[1.5cm]
{\LARGE\bfseries Solución al problema de los Filósofos Comiendo}\\[0.4cm]
{\large Tarea en equipo}\\[1.5cm]
\begin{tabular}{rl}
\textbf{Materia:} & Sistemas Operativos\\
\textbf{Periodo:} & 2026-2\\
\textbf{Profesor:} & Dr. Donald José Rodríguez Úbeda\\
\end{tabular}\\[1.5cm]
{\large\bfseries Integrantes del equipo}\\[0.3cm]
\begin{tabular}{ll}
Angel Alberto Amaya Soria & Expediente 223210350\\
Angel Fernando Bórquez Guerrero & Expediente 219208106\\
Javier Leonardo Miranda Sánchez & Expediente 223215039\\
Oliver Ruiz Beltrán & Expediente 223204115\\
\end{tabular}\\[1.5cm]
Hermosillo, Sonora \hfill Noviembre de 2026
\end{titlepage}

\tableofcontents
\newpage

\section{El problema}
Cinco filósofos están sentados en una mesa redonda. Cada uno pasa la vida alternando entre \emph{pensar} y \emph{comer}. Entre cada par de filósofos hay un palillo (cinco en total), y para comer un filósofo necesita \textbf{los dos palillos} que tiene a sus lados. Cuando termina de comer, los suelta y vuelve a pensar (Dijkstra, 1971; Silberschatz et al., 2018).

\begin{center}
\begin{tikzpicture}[scale=1.0,
  fil/.style={circle,draw,fill=blue!10,minimum size=0.95cm,inner sep=0pt,font=\small},
  pal/.style={rounded corners=1pt,fill=orange!70!black,minimum width=0.12cm,minimum height=0.6cm,inner sep=0pt}]
\draw[gray!40,dotted] (0,0) circle (1.65);
\foreach \i in {0,...,4} {
  \pgfmathsetmacro{\a}{90-72*\i}
  \node[fil] (F\i) at (\a:2.1) {F\i};
  \pgfmathsetmacro{\b}{\a+36}
  \node[pal,rotate=\b-90] at (\b:1.65) {};
  \node[font=\footnotesize,orange!50!black] at (\b:1.2) {p\i};
}
\end{tikzpicture}\\[3pt]
{\footnotesize Figura 1. Cinco filósofos (F0--F4) y cinco palillos (p0--p4). El filósofo $i$ usa los palillos $p_i$ y $p_{(i+1)\bmod 5}$.}
\end{center}

\textbf{¿Por qué es difícil?} Si todos los filósofos se sientan hambrientos al mismo tiempo y cada uno toma primero el palillo de su izquierda, todos quedan con un palillo en la mano esperando el de su derecha, que tiene el vecino. Nadie suelta nada y nadie come: es un \textbf{interbloqueo} (\emph{deadlock}). El problema es un modelo clásico de cómo repartir recursos compartidos (los palillos) entre procesos (los filósofos) sin que se bloqueen entre sí.

\section{La solución: ordenar los palillos}
La idea (orden jerárquico de recursos) es muy sencilla: \textbf{se numeran los palillos de 0 a 4 y todos los filósofos siempre toman primero el palillo de número menor y después el de número mayor.} Los palillos se sueltan al terminar de comer.

\begin{center}
\small
\begin{tabular}{|c|c|c|c|}
\hline
\textbf{Filósofo} & \textbf{Palillos que usa} & \textbf{Toma primero} & \textbf{Toma después}\\
\hline
F0 & p0, p1 & p0 & p1\\ \hline
F1 & p1, p2 & p1 & p2\\ \hline
F2 & p2, p3 & p2 & p3\\ \hline
F3 & p3, p4 & p3 & p4\\ \hline
F4 & p4, p0 & \textbf{p0} & \textbf{p4}\\ \hline
\end{tabular}\\[3pt]
{\footnotesize Tabla 1. Orden de toma de los palillos. F4 es el único que \emph{no} toma primero el de su izquierda.}
\end{center}

Cuatro filósofos toman ``izquierdo y luego derecho'' y el último (F4) lo hace al revés. Ese pequeño cambio rompe el círculo. La figura~2 muestra el peor caso: los cinco tienen hambre a la vez.

\begin{center}
\begin{tikzpicture}[scale=0.95,
  fil/.style={circle,draw,fill=blue!10,minimum size=0.8cm,inner sep=0pt,font=\small},
  esp/.style={-{Stealth[length=2.2mm]},thick,red!70!black}]
% (a) sin orden
\begin{scope}[shift={(-3.6,0)}]
\foreach \i in {0,...,4} \node[fil] (A\i) at (90-72*\i:1.4) {F\i};
\foreach \i/\j in {0/1,1/2,2/3,3/4,4/0} \draw[esp] (A\i) to[bend left=12] (A\j);
\node at (0,-2.35) {\footnotesize (a) Sin orden: hay un ciclo};
\end{scope}
% (b) con orden
\begin{scope}[shift={(3.6,0)}]
\foreach \i in {0,...,4} \node[fil] (B\i) at (90-72*\i:1.4) {F\i};
\foreach \i/\j in {0/1,1/2,2/3,4/0} \draw[esp] (B\i) to[bend left=12] (B\j);
\node[font=\scriptsize,green!40!black,align=center] at (90-72*3-22:2.15) {come};
\node at (0,-2.35) {\footnotesize (b) Con orden: el ciclo se rompe};
\end{scope}
\end{tikzpicture}\\[3pt]
{\footnotesize Figura 2. Una flecha $F_i\to F_j$ significa ``$F_i$ espera un palillo que tiene $F_j$''. En (b), F3 consigue p3 y p4 (nadie más tiene p4) y come; al terminar libera p3 y la cadena avanza.}
\end{center}

\section{Demostración}
Se demuestra que la solución cumple los requisitos de la sección crítica y que no hay \emph{deadlock}. La sección crítica de un filósofo es la parte en que \emph{come}.

\subsection{Exclusión mutua}
\textbf{Requisito:} dos filósofos vecinos nunca comen al mismo tiempo (comparten un palillo).

\textbf{Demostración.} Cada palillo tiene un lock exclusivo: solo un proceso a la vez puede tenerlo. Un filósofo come solo si tiene \emph{sus dos} palillos. Si dos vecinos comieran a la vez, ambos tendrían el palillo que comparten, y eso contradice que el lock sea exclusivo. \hfill$\square$

(Los locks de Linux son \emph{consultivos}: solo protegen si todos los procesos los usan, y aquí todos siguen el mismo protocolo.)

\subsection{Ausencia de deadlock}
Hay interbloqueo solo si se cumplen a la vez cuatro condiciones (Coffman et al., 1971): exclusión mutua, retener y esperar, no apropiación y \textbf{espera circular}. Las tres primeras son inevitables aquí (un palillo no se comparte, se toma uno y se espera el otro, y nadie se lo quita a otro). El ordenamiento elimina la cuarta.

\textbf{Demostración.} Suponga que existe una espera circular $F_1 \to F_2 \to \dots \to F_k \to F_1$, donde cada $F_m$ espera un palillo que tiene $F_{m+1}$. Un filósofo que espera en ese ciclo tiene en la mano su \emph{primer} palillo $a_m$ y espera el segundo $b_m$, con $a_m < b_m$ (siempre toma primero el menor). El palillo que $F_m$ espera lo tiene $F_{m+1}$, quien solo está sosteniendo su primer palillo, o sea $b_m = a_{m+1}$. Entonces
\[
a_1 < b_1 = a_2 < b_2 = a_3 < \dots < a_k < b_k = a_1 ,
\]
lo que implica $a_1 < a_1$, una contradicción. Por lo tanto no puede haber espera circular y no hay \emph{deadlock}. \hfill$\square$

\subsection{Progreso}
\textbf{Requisito:} si hay filósofos con hambre, alguno puede comer en un tiempo finito.

\textbf{Demostración.} Tome a un filósofo hambriento que espera un palillo. Quien lo tiene está comiendo (y terminará en tiempo finito, soltándolo) o está esperando su segundo palillo, que tiene un número \emph{mayor}. Siguiendo la cadena de ``espera a'', los números de palillo crecen cada vez, y solo hay cinco, así que la cadena termina en un filósofo que no espera a nadie: ese puede comer. Así siempre hay alguien que avanza. \hfill$\square$

\subsection{Espera limitada}
La solución garantiza que \textbf{nunca se bloquea el sistema}, pero no garantiza por sí sola el orden de atención: Linux (POSIX) no especifica en qué orden se otorgan los locks a los procesos que esperan, así que, en teoría, un filósofo podría ser rebasado varias veces seguidas (inanición). En la práctica los tiempos de comer son cortos y finitos, y en las pruebas todos los filósofos comieron todas sus rondas. Para tener una garantía estricta habría que añadir una cola de turnos (por ejemplo, un ``mesero'').

\section{Implementación en Linux}
Se simula la mesa con archivos y locks, como en el ejemplo visto en clase de comunicación entre procesos con \cod{fcntl}:
\begin{center}
\small
\begin{tabular}{|>{\raggedright\arraybackslash}p{4.2cm}|>{\raggedright\arraybackslash}p{8.8cm}|}
\hline
\textbf{En el problema} & \textbf{En Linux}\\
\hline
Palillo $i$ & Archivo \cod{palillo$i$.dat}\\ \hline
Tomar el palillo & Lock exclusivo \cod{F\_WRLCK} sobre el archivo\\ \hline
Esperar a que se libere & \cod{fcntl} con \cod{F\_SETLKW} (bloquea al proceso)\\ \hline
Soltar el palillo & \cod{F\_UNLCK}\\ \hline
Filósofo & Un proceso creado con \cod{fork()}\\ \hline
\end{tabular}\\[3pt]
{\footnotesize Tabla 2. Correspondencia entre el problema y la implementación.}
\end{center}
El código completo está en \cod{filosofos.c}. A continuación se explican sus partes.

\subsection{Poner y quitar un lock}
\par\smallskip\noindent\begin{minipage}{\linewidth}
\begin{lstlisting}
void cerrojo(int fd, short tipo) {
   struct flock lock;
   lock.l_type = tipo;      /* F_WRLCK = tomar, F_UNLCK = soltar */
   lock.l_whence = SEEK_SET;
   lock.l_start = 0;
   lock.l_len = 0;          /* 0 = todo el archivo */
   lock.l_pid = getpid();
   while (fcntl(fd, F_SETLKW, &lock) < 0)
      if (errno != EINTR)
         report_and_exit("fcntl falló al manejar el lock...");
}
\end{lstlisting}
\end{minipage}\par\smallskip
Es la misma estructura \cod{flock} del ejemplo del productor y el consumidor, pero con una diferencia importante: se usa \cod{F\_SETLKW} en vez de \cod{F\_SETLK}. Con \cod{F\_SETLK} el proceso fracasa si el palillo está ocupado; con \cod{F\_SETLKW} el proceso \textbf{espera} (se duerme) hasta que el palillo se libera, que es justo lo que necesita un filósofo hambriento.

\subsection{El filósofo}
\par\smallskip\noindent\begin{minipage}{\linewidth}
\begin{lstlisting}
void filosofo(int id, int rondas) {
   int izq = id, der = (id + 1) % N;
   int primero = izq < der ? izq : der;   /* el menor primero */
   int segundo = izq < der ? der : izq;
   int fd1 = abrir_palillo(primero);
   int fd2 = abrir_palillo(segundo);

   for (int r = 1; r <= rondas; r++) {
      usleep(rand() % 300000);            /* piensa */
      cerrojo(fd1, F_WRLCK);              /* toma el 1er palillo */
      usleep(100000);
      cerrojo(fd2, F_WRLCK);              /* toma el 2do palillo */
      /* ... imprime "COME" ... */
      usleep(200000 + rand() % 200000);   /* come (sección crítica) */
      cerrojo(fd2, F_UNLCK);              /* suelta los palillos */
      cerrojo(fd1, F_UNLCK);
   }
}
\end{lstlisting}
\end{minipage}\par\smallskip
Aquí está toda la solución: las dos líneas que calculan \cod{primero} y \cod{segundo} son el \textbf{orden jerárquico}. Para F4, que usa los palillos 4 y 0, el menor es 0, y por eso toma ese primero. El \cod{usleep(100000)} entre un palillo y otro se agregó a propósito: hace más probable un bloqueo si el algoritmo estuviera mal (da tiempo a que todos tengan su primer palillo).

\subsection{El programa principal}
\par\smallskip\noindent\begin{minipage}{\linewidth}
\begin{lstlisting}
for (int i = 0; i < N; i++)             /* crea los archivos-palillo */
   close(abrir_palillo(i));

for (int i = 0; i < N; i++) {           /* un proceso por filósofo */
   pid_t pid = fork();
   if (pid == 0) {
      filosofo(i, rondas);
      exit(0);
   }
}
for (int i = 0; i < N; i++)             /* espera a que todos terminen */
   wait(NULL);
\end{lstlisting}
\end{minipage}\par\smallskip
El padre crea los cinco archivos y lanza cinco procesos hijos (los filósofos). Cada hijo abre sus propios archivos después del \cod{fork}, porque los locks de \cod{fcntl} pertenecen a cada proceso y no se heredan.

\subsection{Compilar y ejecutar}
\par\smallskip\noindent\begin{minipage}{\linewidth}
\begin{lstlisting}[numbers=none,language=bash]
gcc -Wall -o filosofos filosofos.c
./filosofos        # 3 rondas por filósofo (o:  ./filosofos 10)
\end{lstlisting}
\end{minipage}\par\smallskip

\section{Resultados}
Cada línea de salida muestra el tiempo en milisegundos. Un fragmento de una ejecución real:
\par\smallskip\noindent\begin{minipage}{\linewidth}
\begin{lstlisting}[numbers=none,language={},basicstyle=\ttfamily\scriptsize]
[ 1832 ms] Filósofo 3 COME (comida 3)
[ 1856 ms] Filósofo 1 COME (comida 2)
[ 2089 ms] Filósofo 3 TERMINA de comer
[ 2089 ms] Filósofo 4 tomó el palillo 4
[ 2089 ms] Filósofo 4 COME (comida 2)
\end{lstlisting}
\end{minipage}\par\smallskip
Los filósofos 1 y 3 comen a la vez porque no son vecinos. El filósofo 4 solo puede empezar cuando el 3 termina, porque comparten el palillo 4.

\textbf{Verificación.} Se ejecutó el programa \textbf{25} veces y se revisó automáticamente el registro de cada ejecución: ningún par de vecinos comió al mismo tiempo, todos los filósofos completaron sus comidas y ninguna ejecución se bloqueó.

\textbf{Sin el orden.} Al compilar con \cod{-DSIN\_ORDEN} (todos toman primero el palillo izquierdo) aparece el \emph{deadlock} que predice la teoría. Linux lo detecta y, en vez de dejar los procesos dormidos para siempre, hace que \cod{fcntl} falle con el error:
\par\smallskip\noindent\begin{minipage}{\linewidth}
\begin{lstlisting}[numbers=none,language={},basicstyle=\ttfamily\scriptsize]
fcntl falló al manejar el lock...: Resource deadlock avoided
\end{lstlisting}
\end{minipage}\par\smallskip
Es decir, el sistema operativo encuentra el ciclo de espera, el mismo de la figura~2(a), y aborta una de las esperas. Con la solución ordenada ese error nunca ocurre.

\section{Conclusiones}
El problema de los filósofos muestra cómo un conjunto de procesos puede bloquearse al pedir recursos en distinto orden. La solución es simple: \textbf{numerar los recursos y pedirlos siempre en orden creciente}. Así no puede formarse una espera circular, y por tanto no hay deadlock, y la exclusión mutua la da el lock exclusivo de cada palillo. La implementación con archivos y \cod{fcntl} lo comprueba en un sistema real. Como limitación, el orden de atención de los locks no está garantizado por Linux, por lo que la ausencia de inanición no se puede asegurar sin un mecanismo extra de turnos.

\newpage
\section*{Referencias}
\addcontentsline{toc}{section}{Referencias}
\begin{itemize}[leftmargin=1.5em,itemsep=6pt]
\item[] Coffman, E. G., Elphick, M., \& Shoshani, A. (1971). System deadlocks. \emph{ACM Computing Surveys, 3}(2), 67--78.
\item[] Dijkstra, E. W. (1971). Hierarchical ordering of sequential processes. \emph{Acta Informatica, 1}(2), 115--138.
\item[] Kerrisk, M. (2010). \emph{The Linux programming interface}. No Starch Press.
\item[] Silberschatz, A., Galvin, P. B., \& Gagne, G. (2018). \emph{Operating system concepts} (10.ª ed.). Wiley.
\item[] Tanenbaum, A. S., \& Bos, H. (2015). \emph{Modern operating systems} (4.ª ed.). Pearson.
\end{itemize}
\end{document}
